\documentclass[12pt,a4paper]{article}

\usepackage[T1]{fontenc}
\usepackage[utf8]{inputenc}
\usepackage[brazil]{babel}
\usepackage{lmodern}
\usepackage{geometry}
\usepackage{hyperref}
\usepackage{enumitem}
\usepackage{array}

\geometry{margin=2.5cm}
\hypersetup{
  colorlinks=true,
  linkcolor=blue,
  urlcolor=blue
}

\title{Plano de Aula: Controle de Transações}
\date{}

\begin{document}
\maketitle

\noindent
\textbf{Duração:} 2 horas (120 minutos)\\
\textbf{Público-alvo:} Iniciantes\\
\textbf{Recursos Necessários:} Laboratório com computadores, SGBD e cliente SQL (preferencialmente com 2 sessões por dupla), projetor (opcional), quadro e marcadores.

\section{Objetivos da Aula}
\begin{itemize}[leftmargin=*,itemsep=0.25em]
  \item Compreender o conceito de transação e o objetivo de garantir consistência.
  \item Explicar as propriedades ACID (visão conceitual).
  \item Utilizar comandos de controle de transação (BEGIN/COMMIT/ROLLBACK, quando aplicável).
  \item Compreender noções de concorrência e problemas comuns (leituras sujas, perdas de atualização).
\end{itemize}

\section{Metodologia}
Prática orientada por experimentos: as duplas simulam operações que devem ``dar tudo certo ou nada'' (ex.: transferência bancária). Em seguida, executam o mesmo cenário com duas sessões para observar concorrência e importância do isolamento.

\section{Cronograma e Conteúdo Detalhado (120 Minutos)}
\subsection*{Parte 1: Motivação (15 minutos)}
\begin{itemize}[leftmargin=*,itemsep=0.35em]
  \item Cenário: transferência entre contas (debitar e creditar).
  \item Analogia: compra online --- ou paga e confirma, ou não acontece nada.
\end{itemize}

\subsection*{Parte 2: Conceito de Transação e ACID (25 minutos)}
\begin{itemize}[leftmargin=*,itemsep=0.35em]
  \item Atomicidade: tudo ou nada.
  \item Consistência: regras do banco continuam verdadeiras.
  \item Isolamento: transações não ``se atrapalham''.
  \item Durabilidade: após commit, os dados persistem.
\end{itemize}

\subsection*{Parte 3: Comandos de Controle (25 minutos)}
\begin{itemize}[leftmargin=*,itemsep=0.35em]
  \item BEGIN/START TRANSACTION.
  \item COMMIT e ROLLBACK.
  \item SAVEPOINT (quando aplicável) como ``ponto de retorno''.
\end{itemize}

\subsection*{Parte 4: Experimento Guiado (25 minutos)}
\begin{itemize}[leftmargin=*,itemsep=0.35em]
  \item Criar tabela CONTA e inserir saldos iniciais.
  \item Executar uma transferência dentro de transação.
  \item Forçar um erro no meio e aplicar ROLLBACK, observando o estado final.
\end{itemize}

\subsection*{Parte 5: Atividade em Duplas --- Concorrência (25 minutos)}
\begin{itemize}[leftmargin=*,itemsep=0.35em]
  \item Abrir 2 sessões (Sessão A e Sessão B).
  \item Simular duas transações alterando o mesmo saldo.
  \item Observar bloqueios e discutir isolamento (conforme SGBD/configuração).
  \item Entrega: breve relato do que ocorreu e por quê.
\end{itemize}

\subsection*{Parte 6: Encerramento (5 minutos)}
\begin{itemize}[leftmargin=*,itemsep=0.35em]
  \item Revisão: por que transações existem e onde aparecem no mundo real.
  \item Fechamento do curso: ligar modelagem, SQL e integridade com consistência transacional.
\end{itemize}

\end{document}

